\section{The Kuga--Satake construction}\label{sec:ks}

\Cref{thm:hyperplane} blocks induction on dimension through hyperplane
sections in the primitive middle degree. One classical construction does reach
a middle degree Hodge class without any such induction, and we set it out in
full, both because it shows what a successful route looks like and because
its boundary is numerical and can be located. It is the construction of Kuga
and Satake \cite{KS67}, which builds on Satake's Clifford algebra families and
their Riemann forms \cite{Sat66}. We follow the form given by Deligne
\cite{Del72}, to whom the Hodge-theoretic properties that make it useful are
due, in particular the embedding of \Cref{prop:ksembed}. After building it we
show that \Cref{thm:hyperplane} has no bearing on it. For abelian varieties of
Weil type with Hodge group $\SU(V,H)$ we then determine when it applies to a
piece of $H^{2}$: exactly when $n=1$, where the Weil line consists of divisor
classes and there is nothing to prove, or $n=2$ with trivial discriminant,
where the Weil line lies in $H^{4}$.

\subsection{The construction}

\begin{setupenv}[K3 type]\label{setup:k3type}
Let $T$ be a $\QQ$-vector space of dimension $b$ with a nondegenerate
symmetric form $q$, and let $T$ carry a Hodge structure of weight two,
$T_{\CC}=T^{2,0}\oplus T^{1,1}\oplus T^{0,2}$, polarised by $q$, with
\[
  \dim_{\CC}T^{2,0}=1 .
\]
Such a structure is said to be \emph{of K3 type}. Write $T^{2,0}=\CC\sigma$.
The Hodge--Riemann relations give $q(\sigma,\sigma)=0$ and
$q(\sigma,\bbar{\sigma})>0$, and make $q$ negative definite on the real points
of $T^{1,1}$; so the real plane
$P=\langle\Re\sigma,\Im\sigma\rangle\subset T_{\RR}$ is positive definite and
$q$ has signature $(2,b-2)$ on $T_{\RR}$. Since $q(\Re\sigma,\Im\sigma)=0$
and $q(\Re\sigma)=q(\Im\sigma)$, we may normalise $\sigma$ so that
$e_{1}=\Re\sigma$ and $e_{2}=\Im\sigma$ are orthonormal.
\end{setupenv}

\begin{definition}[Clifford algebra]\label{def:clifford}
$C(T,q)$ is the quotient of the tensor algebra on $T$ by the ideal generated
by $v\otimes v-q(v)$ for $v\in T$. It is $\ZZ/2$-graded, $C=C^{+}\oplus C^{-}$,
by the parity of the word length.
\end{definition}

\begin{lemma}\label{lem:cliffdim}
For an orthogonal basis $f_{1},\dots,f_{b}$ of $T$ the products $f_{S}$ over
subsets $S\subseteq\{1,\dots,b\}$, taken in increasing order, form a basis of
$C(T,q)$. Hence $\dim_{\QQ}C=2^{b}$ and $\dim_{\QQ}C^{\pm}=2^{b-1}$, and
$C^{-}\cdot C^{+}\cdot C^{-}\subseteq C^{+}$.
\end{lemma}

\begin{proof}
The relations $f_{i}^{2}=q(f_{i})$ and $f_{i}f_{j}=-f_{j}f_{i}$ for $i\ne j$,
the latter obtained by expanding $(f_{i}+f_{j})^{2}=q(f_{i})+q(f_{j})$, let one
rewrite any word in the $f_{i}$ as a scalar times an increasing product, so the
$f_{S}$ span; they are independent by the standard filtration argument
comparing $C$ with $\bigwedge^{*}T$. The dimension counts are
$\sum_{k}\binom{b}{k}=2^{b}$ and
$\sum_{k\ \mathrm{even}}\binom{b}{k}=2^{b-1}$. The last inclusion holds
because the grading is multiplicative and odd plus even plus odd is even.
\end{proof}

\begin{lemma}[The complex structure]\label{lem:cliffJ}
With $e_{1},e_{2}$ as in \Cref{setup:k3type}, put $J=e_{1}e_{2}\in C^{+}\otimes_{\QQ}\RR$.
Then $J^{2}=-1$, and left multiplication by $J$ is an $\RR$-linear operator on
$C^{+}\otimes_{\QQ}\RR$ whose square is $-\id$.
\end{lemma}

\begin{proof}
$J^{2}=e_{1}e_{2}e_{1}e_{2}=-e_{1}e_{1}e_{2}e_{2}=-q(e_{1})q(e_{2})=-1$,
since the orthogonal vectors $e_{1}$ and $e_{2}$ anticommute and
$e_{i}^{2}=q(e_{i})=1$.
Left multiplication by an element of $C^{+}$ preserves $C^{+}$, and the square
of left multiplication by $J$ is left multiplication by $J^{2}=-1$.
\end{proof}

\Cref{fig:ks} draws the positive plane $P$, the null cone of $q$ and the
period domain in the smallest case $b=3$.

\begin{definition}[Kuga--Satake variety]\label{def:ksvariety}
Fix a lattice $C^{+}_{\ZZ}\subset C^{+}$ stable under multiplication. The
complex torus
\[
  \mathrm{KS}(T)=\bigl(C^{+}\otimes_{\QQ}\RR,\,J\bigr)\big/\,C^{+}_{\ZZ}
\]
has dimension $\tfrac12\dim_{\QQ}C^{+}=2^{b-2}$. Satake's construction of a
Riemann form on $C^{+}$ compatible with $J$ \cite{Sat66} makes it an abelian
variety, and by construction
\[
  H^{1}\bigl(\mathrm{KS}(T),\QQ\bigr)=C^{+}(T)
\]
as a Hodge structure of weight one, the Hodge decomposition being the
eigenspace decomposition of left multiplication by $J$.
\end{definition}

\begin{example}\label{ex:ksdims}
For a projective K3 surface, $T$ is the transcendental lattice inside
$H^{2}(X,\QQ)$; taking all of $H^{2}$ with $b=22$ gives
$\dim\mathrm{KS}=2^{20}=1{,}048{,}576$. At the other end, $b=3$ gives
$\dim\mathrm{KS}=2$, an abelian surface, and $b=6$ gives
$\dim\mathrm{KS}=16$; these are the two cases that occur below, for $n=1$ and
for $n=2$.
\end{example}

\subsection{The transport of Hodge classes}

\begin{proposition}[The embedding]\label{prop:ksembed}
Fix $v_{0}\in T$ with $q(v_{0})\ne0$. The map
\[
  \Psi\colon T\longrightarrow\End_{\QQ}\bigl(C^{+}\bigr),
  \qquad
  \Psi(v)(x)=v\,x\,v_{0},
\]
is injective and is a morphism of Hodge structures
$T(1)\to\End\bigl(H^{1}(\mathrm{KS}(T),\QQ)\bigr)$, both sides having weight
zero. Consequently a Hodge class in $T$, that is a rational class of type
$(1,1)$, is carried to a Hodge class of type $(0,0)$ in
$\End(H^{1}(\mathrm{KS}(T),\QQ))$, that is to an element of
$\End^{0}(\mathrm{KS}(T))$.
\end{proposition}

\begin{proof}
The map lands in $\End(C^{+})$ by the last part of \Cref{lem:cliffdim}, since
$v,v_{0}\in T\subset C^{-}$. It is injective because $\Psi(v)(1)=v\,v_{0}$
and $v_{0}$ is invertible in $C(T,q)$, with $v_{0}^{-1}=v_{0}/q(v_{0})$, so
$\Psi(v)=0$ forces $v=0$. That it respects the Hodge structures is
Deligne's computation \cite[\S4]{Del72}: the Hodge structure on
$\End(C^{+})$ is induced by $J$ acting on both sides, and $\Psi$ intertwines
the weight-two structure on $T$, twisted once, with it. The last sentence
follows from the definition of a morphism of Hodge structures and the
Lefschetz theorem on $(1,1)$ classes, in the form which says that for an
abelian variety $B$ a Hodge class of type $(0,0)$ in $\End(H^{1}(B,\QQ))$ is
an element of $\End^{0}(B)$.
\end{proof}

\begin{remark}[Algebraicity of the transport]%
\label{rem:kstransport}
An element of $\End^{0}(\mathrm{KS}(T))$ is a rational multiple of an
endomorphism, whose graph is an algebraic cycle on
$\mathrm{KS}(T)\times\mathrm{KS}(T)$, so the image of a Hodge class under
$\Psi$ is algebraic. Carrying that conclusion back to the variety $X$ from
which $T$ came requires the correspondence $\Psi$ itself to be induced by an
algebraic cycle on $X\times\mathrm{KS}(T)\times\mathrm{KS}(T)$. What is
proved in general is that $\Psi$ is absolute Hodge, by
Deligne's theorem for abelian varieties together with the analysis of the
Kuga--Satake correspondence, and that it is motivated in Andr\'e's sense
\cite{And96}; van Geemen \cite{vG00} surveys the cases in which algebraicity
has been established. The Kuga--Satake construction therefore converts the
Hodge conjecture for $T$ into the algebraicity of one specific
correspondence.
\end{remark}

\subsection{Independence from the hyperplane obstruction}

\begin{proposition}\label{prop:ksnotrestriction}
\Cref{thm:hyperplane,thm:returntrip} do not constrain the Kuga--Satake route.
\end{proposition}

\begin{proof}
Both theorems are statements about the pair of maps $j^{*}$ and $j_{*}$
attached to the inclusion $j\colon Y\hookrightarrow X$ of a smooth member of an
ample linear system, and both proofs use two properties of that pair which a
general correspondence does not have: the Lefschetz hyperplane theorem, which
makes $j^{*}$ an isomorphism in degrees at most $d-2$ and, through
\Cref{lem:gysinrestrict}, forces the vanishing of \Cref{thm:hyperplane} in
the primitive middle degree; and the identity $j_{*}j^{*}=\,\cdot\cup L$ of
\Cref{lem:gysinrestrict}, which makes the return trip multiply by the
polarisation instead of recovering the class. The Kuga--Satake
correspondence is a cohomology class on a triple product of varieties of
unrelated dimensions, $\dim\mathrm{KS}(T)=2^{b-2}$ against $\dim X$, and it has
neither property; in particular the composite of $\Psi$ with its transpose is
a nonzero multiple of the identity of $T$, not a power of a polarisation of
$X$. So
neither of the two properties on which \Cref{thm:hyperplane} rests holds for
$\Psi$.
\end{proof}

\Cref{prop:ksnotrestriction} shows what an escape from \Cref{thm:hyperplane}
must look like: a correspondence whose composite with its transpose is not the
Lefschetz operator. \Cref{sec:escape} makes the analogous identification for
\Cref{thm:divisor,thm:character}.

\subsection{Abelian varieties of Weil type}

The limitation of the Kuga--Satake construction is the hypothesis
$\dim T^{2,0}=1$. For abelian varieties of Weil type the rational sub-Hodge
structures of $H^{2}$, and hence the possible values of $\dim T^{2,0}$, can be
listed completely.

\begin{proposition}[The pieces of \texorpdfstring{$H^{2}$}{H2}]\label{prop:h2pieces}
Let $(A,\iota,\eta)$ be of $(K,-1,n)$-Weil type. In the character grading of
\Cref{prop:grading},
\begin{equation}\label{eq:h2split}
  H^{2}(A,\CC)
  =\textstyle\bigwedge^{2}V_{+}\;\oplus\;\bigl(V_{+}\otimes V_{-}\bigr)
   \;\oplus\;\bigwedge^{2}V_{-},
\end{equation}
the middle summand is the isotypic component of the norm character and is
defined over $\QQ$, and the outer two are interchanged by conjugation so their
sum is defined over $\QQ$. The Hodge numbers are
\[
  \bigl(h^{2,0},h^{1,1},h^{0,2}\bigr)=
  \begin{cases}
    \bigl(\tbinom{n}{2},\,n^{2},\,\tbinom{n}{2}\bigr)
      & \text{on }\bigwedge^{2}V_{\pm},\\[4pt]
    \bigl(n^{2},\,2n^{2},\,n^{2}\bigr)
      & \text{on }V_{+}\otimes V_{-},
  \end{cases}
\]
and the total is $2\tbinom{2n}{2}+4n^{2}=\tbinom{4n}{2}=\dim H^{2}(A,\QQ)$.
\end{proposition}

\begin{proof}
The splitting is \eqref{eq:grading} at $k=2$. The characters are
$\tau^{2}$, $\tau\bbar{\tau}=\Nm(\tau)$ and $\bbar{\tau}^{2}$, so the middle
summand is the norm-isotypic component, and it is self-conjugate, hence defined
over $\QQ$; the outer two are exchanged. For the Hodge numbers use
\eqref{eq:bisplit} and \Cref{lem:balanced}: $V_{\pm}^{1,0}$ and
$V_{\pm}^{0,1}$ have dimension $n$ each, so
$(\bigwedge^{2}V_{+})^{2,0}=\bigwedge^{2}V^{1,0}_{+}$ has dimension
$\binom{n}{2}$ and $(\bigwedge^{2}V_{+})^{1,1}=V^{1,0}_{+}\otimes V^{0,1}_{+}$
has dimension $n^{2}$, while
$(V_{+}\otimes V_{-})^{2,0}=V^{1,0}_{+}\otimes V^{1,0}_{-}$ has dimension
$n^{2}$ and the $(1,1)$ part has dimension $2n^{2}$. The total is
$2(2\binom{n}{2}+n^{2})+4n^{2}=2(2n^{2}-n)+4n^{2}=8n^{2}-2n=\binom{4n}{2}$.
\end{proof}

\begin{proposition}[Sub-Hodge structures of K3 type]%
\label{prop:normpart}
Let $(A,\iota,\eta)$ be of $(K,-1,n)$-Weil type with $\Hg(A)=\SU(V,H)$. Then
$H^{2}(A,\QQ)$ is the direct sum of the three rational sub-Hodge structures
\[
  \QQ\,\eta,
  \qquad
  \bigl(V_{+}\otimes V_{-}\bigr)_{\QQ}\ominus\QQ\eta,
  \qquad
  R=\Bigl(\textstyle\bigwedge^{2}V_{+}\oplus\bigwedge^{2}V_{-}\Bigr)_{\QQ},
\]
with $h^{2,0}$ equal to $0$, $n^{2}$ and $n(n-1)$ respectively. The first two
are irreducible, and so is $R$ when $n\ge3$. When $n=2$, the structure $R$ is
irreducible if the discriminant of $(A,\iota,\eta)$ in the sense of
\Cref{def:disc} is nontrivial, and it is the direct sum of two isomorphic
sub-Hodge structures of K3 type, each of dimension six, if the discriminant
is trivial. Consequently $H^{2}(A,\QQ)$ contains a rational sub-Hodge
structure of K3 type if and only if either $n=1$, or $n=2$ and the
discriminant is trivial.
\end{proposition}

\begin{proof}
A rational sub-Hodge structure $W\subseteq H^{2}(A,\QQ)$ is stable under
$\Hg(A)$: the orthogonal projector onto $W$ for the polarisation is a Hodge
class in $\End(H^{2})$, hence is fixed by $\Hg(A)$, hence $W$ is a
$\Hg(A)$-submodule. Conversely a $\Hg(A)$-submodule is a sub-Hodge structure,
because the Hodge decomposition is the eigenspace decomposition of an element
of $\Hg(A)(\RR)$. So the rational sub-Hodge structures are the
$\QQ$-subrepresentations of $\SU(V,H)$. For the same reason the Hodge numbers
of a sub-Hodge structure depend only on its complexification as a
representation of $\Hg(A)_{\CC}$.

After complexification $\SU(V,H)$ becomes $\SL(V_{+})$, acting on $V_{+}$
tautologically and on $V_{-}\cong V_{+}^{\vee}$ dually. In \eqref{eq:h2split}
the outer summands $\bigwedge^{2}V_{+}$ and
$\bigwedge^{2}V_{-}\cong\bigwedge^{2}V_{+}^{\vee}$ are irreducible, and the
middle summand is $V_{+}\otimes V_{+}^{\vee}=\End(V_{+})
=\mathfrak{sl}(V_{+})\oplus\CC$, the sum of the adjoint representation and the
trivial one. Conjugation fixes the two pieces of the middle summand and
exchanges the outer two, so the three listed subspaces are rational sub-Hodge
structures, and the first two, having irreducible complexifications, are
irreducible. The trivial summand is spanned by $\eta$ by \Cref{lem:NSchar},
and has $h^{2,0}=0$ since $\eta$ is of type $(1,1)$. The adjoint summand has
$h^{2,0}=n^{2}$, since removing the one-dimensional trivial piece from
$V_{+}\otimes V_{-}$ removes a $(1,1)$ class only. The structure $R$ has
$h^{2,0}=2\binom{n}{2}=n(n-1)$ by \Cref{prop:h2pieces}.

The complexification of a rational subrepresentation of $R$ is a
subrepresentation of $\bigwedge^{2}V_{+}\oplus\bigwedge^{2}V_{-}$ stable under
conjugation. For $n\ge3$ the two summands are not isomorphic, because
$\bigwedge^{2}V_{+}^{\vee}\cong\bigwedge^{2n-2}V_{+}$ and $2n-2\ne2$. The
subrepresentations are then $0$, the two summands and the whole; conjugation
exchanges the summands, so only $0$ and the whole are stable under it, and
$R$ is irreducible.
For $n=1$ both summands are the trivial representation, and $R$ is the plane
of Hodge classes $\HW(A)\subset H^{2}(A,\QQ)$.

Let $n=2$. Then $\bigwedge^{2}V_{+}^{\vee}\cong\bigwedge^{2}V_{+}\otimes
(\bigwedge^{4}V_{+})^{\vee}\cong\bigwedge^{2}V_{+}$ as representations of
$\SL(V_{+})$, so $\End_{\Hg(A)}(R)\otimes\CC\cong M_{2}(\CC)$ and
$\End_{\Hg(A)}(R)$ is a quaternion algebra over $\QQ$. To identify it, put
$V=H^{1}(A,\QQ)$, a $K$-vector space of dimension four by \Cref{lem:Kspace}.
The natural map $\bigwedge^{2}_{\QQ}V\to\bigwedge^{2}_{K}V$ is
$\SU(V,H)$-equivariant, and its kernel has complexification
$V_{+}\otimes V_{-}$, so it restricts to an isomorphism of $R$ onto
$\bigwedge^{2}_{K}V$; in particular $K\subseteq\End_{\Hg(A)}(R)$. The wedge
product $B\colon\bigwedge^{2}_{K}V\times\bigwedge^{2}_{K}V\to
\bigwedge^{4}_{K}V\cong K$ is a symmetric $K$-bilinear form, invariant under
$\SU(V,H)$ because $\SU(V,H)$ acts trivially on $\bigwedge^{4}_{K}V$. Let
$H_{2}$ be the hermitian form induced by $H$ on $\bigwedge^{2}_{K}V$, and let
$\varphi$ be the $K$-antilinear map defined by
$B(x,y)=H_{2}\bigl(y,\varphi(x)\bigr)$. Then $\varphi$ commutes with
$\SU(V,H)$, so $\varphi^{2}$ is a $K$-linear endomorphism commuting with
$\SU(V,H)$; it is a scalar by Schur's lemma, and a rational one because it
commutes with $\varphi$. Choose a $K$-basis $e_{1},\dots,e_{4}$ of $V$
orthogonal for $H$, with $H(e_{i},e_{i})=h_{i}\in\QQ^{\times}$, and identify
$\bigwedge^{4}_{K}V$ with $K$ through $e_{1}\wedge e_{2}\wedge e_{3}\wedge
e_{4}$. Then $\varphi(e_{1}\wedge e_{2})=(h_{3}h_{4})^{-1}e_{3}\wedge e_{4}$
and $\varphi(e_{3}\wedge e_{4})=(h_{1}h_{2})^{-1}e_{1}\wedge e_{2}$, so
$\varphi^{2}=(\det H)^{-1}$. Hence $\End_{\Hg(A)}(R)=K\oplus K\varphi$ is the
quaternion algebra $(-d,\det H)_{\QQ}$. It is split if and only if
$\det H\in\Nm_{K/\QQ}(\Kx)$, that is, since $(-1)^{2}\det H=\det H$, if and
only if the discriminant is trivial. If it is not split it is a division
algebra, and $R$ is irreducible. If it is split, $R=R_{1}\oplus R_{2}$ with
$R_{1}\cong R_{2}$ of dimension six, and the complexification of each $R_{i}$
is isomorphic to $\bigwedge^{2}V_{+}$ as a representation of $\SL(V_{+})$, so
by \Cref{prop:h2pieces} each $R_{i}$ has Hodge numbers $(1,4,1)$. Each $R_{i}$
lies in the primitive part of $H^{2}(A,\QQ)$, which is the $\Hg(A)$-stable
complement of $\QQ\eta$, and is polarised by the restriction of the
Hodge--Riemann form; so $R_{i}$ is of K3 type.

Finally, a rational sub-Hodge structure of $H^{2}(A,\QQ)$ is a direct sum of
irreducible ones, each isomorphic to $\QQ\eta$, to the adjoint structure, to
$R$, or, for $n=2$ and trivial discriminant, to $R_{1}$; and $h^{2,0}$ is
additive. A structure of K3 type has $h^{2,0}=1$. For $n=1$ the adjoint
structure has $h^{2,0}=1$. For $n\ge3$ the nonzero values $n^{2}$ and
$n(n-1)$ both exceed $1$. For $n=2$ the values are $0$, $4$, and $2$ for $R$
or $1$ for $R_{1}$, so a structure of K3 type exists if and only if the
discriminant is trivial.
\end{proof}

\begin{corollary}[Where the construction applies]\label{cor:ksboundary}
Let $(A,\iota,\eta)$ be of $(K,-1,n)$-Weil type with $\Hg(A)=\SU(V,H)$. The
Kuga--Satake construction applies to a rational sub-Hodge structure of
$H^{2}(A,\QQ)$ if and only if either $n=1$, or $n=2$ and the discriminant is
trivial. For $n=1$ the Weil line lies in $H^{2}(A,\QQ)$ and consists of
divisor classes, hence of algebraic classes by the Lefschetz theorem on
$(1,1)$ classes, so nothing needs transporting. For $n=2$ the Weil line lies
in $H^{4}(A,\QQ)$, outside the structures to which the construction applies.
\end{corollary}

\begin{proof}
The first clause is \Cref{prop:normpart} with \Cref{setup:k3type}. For $n=1$
we have $2n=2$, so $\HW(A)\subseteq H^{2}(A,\QQ)\cap H^{1,1}(A)$, which is
$\NS(A)\otimes\QQ$ by the Lefschetz theorem on $(1,1)$ classes. For $n=2$ the
Weil line lies in degree $2n=4$.
\end{proof}

\begin{remark}[Higher weight]\label{rem:nohigherks}
For $n\ge2$ the Weil line sits in $H^{2n}(A,\QQ)$ with $2n\ge4$, and the
Clifford construction has no known counterpart in weight greater than two. It
turns a quadratic form into an algebra; the polarisation of a Hodge structure
of weight $w$ is symmetric only when $w$ is even, and even then there is no
reason for its Clifford algebra to carry the right complex structure. One
can instead apply the construction to a sub-Hodge structure of $H^{2}$ and
try to recover the degree $2n$ class from the resulting abelian variety.
\Cref{prop:normpart} rules this out for $n\ge3$, and for $n=2$ at nontrivial
discriminant. At $n=2$ and trivial discriminant the two structures of K3 type
exist, but the Hodge conjecture for abelian fourfolds of Weil type is in any
case a theorem of Markman \cite{Mar25a,Mar25b}, proved at trivial
discriminant also by Floccari and Fu \cite{FF26}. What remains conceivable for
$n\ge3$ is a construction of Kuga--Satake shape for the weight $2n$ structure
directly, producing an auxiliary variety on which the Weil line becomes a
divisor class or an endomorphism. None is known, and \Cref{thm:character} does
not obstruct one: such a construction is a correspondence and need not
produce a $\Kx$-eigenclass. This is the gap of \Cref{rem:nonequivariant},
seen from the other side.
\end{remark}

\begin{figure}[htbp]
\centering
  \includegraphics[width=0.88\textwidth]{fig_lightcone}
\caption{The geometry from which the Kuga--Satake construction starts, in the
smallest case $b=3$. The quadratic form is $q=x^{2}+y^{2}-z^{2}$, of
signature $(2,1)$, with its null cone $q=0$; the horizontal plane
$P=\{z=0\}$ is positive definite and meets the cone only at the origin. The
construction needs such a plane, and \Cref{setup:k3type} supplies one:
$P=\langle\Re\sigma,\Im\sigma\rangle$ for the period $\sigma$ spanning
$T^{2,0}$. An orthonormal basis $e_{1},e_{2}$ of $P$ gives $J=e_{1}e_{2}$ with
$J^{2}=-1$ by \Cref{lem:cliffJ}, and $J$ is the complex structure of
$\mathrm{KS}(T)$. A positive definite plane $P$ is determined by the negative
line $P^{\perp}$, which meets the upper sheet of the hyperboloid $q=-1$ in one
point $w_{P}$; that sheet is the period domain, and a second plane $P'$ with
its point $w_{P'}$ is drawn for comparison. The hypothesis behind all of this
is $\dim T^{2,0}=1$, and by \Cref{prop:normpart} an abelian variety of Weil
type with Hodge group $\SU(V,H)$ supplies it in $H^{2}$ only when $n=1$, or
$n=2$ at trivial discriminant.}
\label{fig:ks}
\end{figure}
